\documentclass[11pt,a4paper]{article}
\usepackage{amsmath,amssymb}
\usepackage[margin=2.5cm]{geometry}
\usepackage{array}
\usepackage{bm}

\title{Self-Attention and Cross-Attention via Mamba-3, and Their Application to Images}
\author{Masahiro Ogawa}
\date{\today}

\begin{document}
\maketitle

\begin{abstract}
This note derives, step by step, how the Mamba-3 state space model
(SSM) can be used as a drop-in replacement for both self-attention and
cross-attention, and how the resulting attention blocks are assembled
into an image model. We follow the notation of the original Mamba-3
paper so every equation can be checked against the source. The
guiding principle is that every algebraic move is shown with its
justification. A reader who knows matrix multiplication and
elementary calculus should be able to follow the whole note without
guessing.
\end{abstract}

\tableofcontents

%% ================================================================
\section{Preliminaries and Notation}
%% ================================================================

We summarize the symbols used throughout. The names match
\cite{mamba3}, Section~2 and Section~3.

\subsection{Sizes}

\begin{itemize}
  \item $T$: sequence length (number of tokens). For images we will
        set $T = L = HW/P^2$ where $H,W$ are the image height and
        width and $P$ is the patch size.
  \item $D$: model (channel) dimension.
  \item $N$: SSM state dimension (per SSM head).
  \item $R$: MIMO rank (Mamba-3, \S 3.3). For SISO, $R=1$.
\end{itemize}

\subsection{Tokens and projections}

\begin{itemize}
  \item $\bm{X} \in \mathbb{R}^{T\times D}$ is the input sequence;
        the $t$-th row $\bm{x}_t \in \mathbb{R}^{D}$ is one token.
  \item $\bm{B}, \bm{C} \in \mathbb{R}^{T\times N}$ are data-dependent
        projections; the $t$-th rows are $\bm{B}_t, \bm{C}_t \in
        \mathbb{R}^{N}$.
  \item $\Delta_t \in \mathbb{R}_{>0}$ is a data-dependent step size.
  \item $A_t \in \mathbb{R}_{<0}$ is a data-dependent (in Mamba-3) or
        data-independent (in Mamba-2) scalar that governs decay.
  \item $\bm{h}_t \in \mathbb{R}^{N}$ is the hidden state at time $t$
        (SISO; see \S\ref{sec:mimo} for MIMO).
  \item $\bm{Y} \in \mathbb{R}^{T\times D}$ is the output sequence.
\end{itemize}

In terms of an input token $\bm{x}_t$, the learned projections are
\begin{align}
  \bm{B}_t       &= \bm{W}_B\,\bm{x}_t  & \bm{W}_B &\in \mathbb{R}^{N\times D}, \label{eq:projB}\\
  \bm{C}_t       &= \bm{W}_C\,\bm{x}_t  & \bm{W}_C &\in \mathbb{R}^{N\times D}, \label{eq:projC}\\
  \Delta_t       &= \mathrm{Softplus}(\bm{w}_\Delta^\top \bm{x}_t)
                                         & \bm{w}_\Delta &\in \mathbb{R}^{D}, \label{eq:projD}\\
  A_t            &= -\mathrm{Softplus}(\bm{w}_A^\top \bm{x}_t)
                                         & \bm{w}_A &\in \mathbb{R}^{D} \text{ (optional data-dep.)}. \label{eq:projA}
\end{align}
Here $\mathrm{Softplus}$ is the smooth, strictly positive activation
\begin{equation}\label{eq:softplus}
  \mathrm{Softplus}(z) := \log\bigl(1 + e^{z}\bigr),
  \qquad \mathrm{Softplus}:\mathbb{R}\to\mathbb{R}_{>0},
\end{equation}
so \eqref{eq:projD} guarantees $\Delta_t > 0$ and \eqref{eq:projA}
guarantees $A_t < 0$, which is what the discretization in
\eqref{eq:abc} requires for $\alpha_t = e^{\Delta_t A_t} \in (0,1)$.

The value projection that we will feed into the SSM is
\begin{equation}\label{eq:projV}
  \bm{v}_t = \mathrm{SiLU}(\bm{W}_V\,\bm{x}_t),
  \qquad \bm{W}_V \in \mathbb{R}^{D\times D},
\end{equation}
where $\mathrm{SiLU}$ (a.k.a.\ Swish) is the smooth, self-gated activation
\begin{equation}\label{eq:silu}
  \mathrm{SiLU}(z) := z\,\sigma(z) = \frac{z}{1 + e^{-z}},
  \qquad
  \sigma(z) := \frac{1}{1 + e^{-z}},
\end{equation}
applied element-wise to the vector $\bm{W}_V\bm{x}_t \in \mathbb{R}^{D}$.

\subsection{RMSNorm and BCNorm}
\label{sec:bcnorm}

Mamba-3 normalizes the key/query projections $\bm{B}_t,\bm{C}_t$ before
using them in the SSD kernel. The underlying primitive is
\emph{RMSNorm} \cite{rmsnorm}, which rescales a vector by its
root-mean-square (no mean subtraction, unlike LayerNorm):
\begin{equation}\label{eq:rmsnorm}
  \mathrm{RMSNorm}(\bm{u}) := \frac{\bm{u}}{\sqrt{\tfrac{1}{N}\sum_{i=1}^{N} u_i^{2} + \varepsilon}}\;\odot\;\bm{g},
  \qquad \bm{u}\in\mathbb{R}^{N},
\end{equation}
where $\bm{g}\in\mathbb{R}^{N}$ is a learnable per-channel gain,
$\varepsilon>0$ is a small constant, and $\odot$ is the Hadamard
product. The output has unit RMS up to the learned gain, which
stabilizes the inner products $\bm{C}_i^\top\bm{B}_j$ that appear in
\eqref{eq:ssd}.

\emph{BCNorm} (\cite{mamba3} \S 3.4) is simply RMSNorm applied
independently to $\bm{B}_t$ and $\bm{C}_t$, followed by a learnable
channel-wise bias:
\begin{equation}\label{eq:bcnorm}
  \bm{B}_t \leftarrow \mathrm{RMSNorm}(\bm{B}_t) + \bm{b}_B,
  \qquad
  \bm{C}_t \leftarrow \mathrm{RMSNorm}(\bm{C}_t) + \bm{b}_C,
  \qquad \bm{b}_B,\bm{b}_C \in \mathbb{R}^{N}.
\end{equation}
We will write ``apply BCNorm'' as shorthand for the substitution
\eqref{eq:bcnorm}; each occurrence of $\bm{B}_t,\bm{C}_t$ in later
equations is understood to be the post-BCNorm version when the
algorithm calls for it.

\subsection{Recurrence, parallel form, and decay factors}

Mamba-3 uses the \emph{exponential-trapezoidal} discrete recurrence of
\cite{mamba3} eq.~(5)–(6):
\begin{equation}\label{eq:rec3}
  \bm{h}_t = \alpha_t\,\bm{h}_{t-1}
            + \beta_t\,\bm{B}_{t-1}\bm{x}_{t-1}
            + \gamma_t\,\bm{B}_t\bm{x}_t,
  \qquad
  \bm{y}_t = \bm{C}_t^\top\bm{h}_t,
\end{equation}
with the three data-dependent scalars
\begin{equation}\label{eq:abc}
  \alpha_t := e^{\Delta_t A_t},
  \qquad
  \beta_t  := (1 - \lambda_t)\,\Delta_t\,e^{\Delta_t A_t},
  \qquad
  \gamma_t := \lambda_t\,\Delta_t,
\end{equation}
where $\lambda_t \in [0,1]$ is a learned, data-dependent convex weight.
Setting $\lambda_t \equiv 1$ recovers Mamba-2's two-term recurrence
\begin{equation}\label{eq:rec2}
  \bm{h}_t = \alpha_t\,\bm{h}_{t-1} + \gamma_t\,\bm{B}_t\,\bm{x}_t,
  \qquad
  \bm{y}_t = \bm{C}_t^\top\,\bm{h}_t,
\end{equation}
which is \cite{mamba3} eq.~(1). Mamba-3's \emph{parallel} form
(\cite{mamba3} eq.~(2)) reads
\begin{equation}\label{eq:ssd}
  \bm{Y} = \bigl(\bm{L}\odot\bm{C}\bm{B}^\top\bigr)\bm{X},
\end{equation}
with $\odot$ the Hadamard product and $\bm{L}\in\mathbb{R}^{T\times T}$
a structured lower-triangular mask that we derive explicitly in
Section~\ref{sec:ssd}. This single equation is the bridge between
recurrences and attention: the next section is devoted to explaining
exactly why.

\subsection{What this note actually does}

Using only the objects above, we will:

\begin{enumerate}
  \item Derive \eqref{eq:ssd} from \eqref{eq:rec2} and
        \eqref{eq:rec3}, and read it as \emph{linear attention with a
        learned mask} (Section~\ref{sec:ssd}).
  \item Plug real projections into \eqref{eq:ssd} to get a Mamba-3
        self-attention block (Section~\ref{sec:self}).
  \item Repeat for two different sequences and obtain Mamba-3 cross-
        attention, in two flavours: state-compressed (cheap,
        constant-memory) and token-level (exact SSD analogue of
        softmax cross-attention) (Section~\ref{sec:cross}).
  \item Patchify an image, make the scan bidirectional, and use the
        self/cross-attention blocks to build an image encoder; then
        combine several views via cross-view attention for the 3-D
        setting (Section~\ref{sec:images}).
  \item Derive the \emph{non-causal SSD} (NC-SSD) of \cite{vssd}
        as an algebraic alternative to bidirectional scans, and show
        that it collapses the structured $T\times T$ mask $\bm{L}$ to
        a single learned $T$-vector $\bm{m}$
        (Section~\ref{sec:ncssd}).
\end{enumerate}

%% ================================================================
\section{Why SSD $\equiv$ linear attention (the duality, derived)}
\label{sec:ssd}
%% ================================================================

We derive \eqref{eq:ssd} directly from the recurrence. Work entirely
with the two-term form \eqref{eq:rec2} first; the three-term form
\eqref{eq:rec3} is a corollary (\S\ref{sec:trap}).

\subsection{Unrolling the two-term recurrence}

Starting from $\bm{h}_0 = \bm{0}$ and applying \eqref{eq:rec2} once:
\begin{equation}\label{eq:un1}
  \bm{h}_1 = \alpha_1\,\bm{h}_0 + \gamma_1\,\bm{B}_1\bm{x}_1
           = \gamma_1\,\bm{B}_1\bm{x}_1.
\end{equation}
Apply it once more:
\begin{align}
  \bm{h}_2 &= \alpha_2\,\bm{h}_1 + \gamma_2\,\bm{B}_2\bm{x}_2 \nonumber \\
           &= \alpha_2\,\gamma_1\,\bm{B}_1\bm{x}_1
             + \gamma_2\,\bm{B}_2\bm{x}_2. \label{eq:un2}
\end{align}
By induction on $t$ we claim
\begin{equation}\label{eq:un-claim}
  \bm{h}_t \;=\; \sum_{j=1}^{t}
     \underbrace{\Bigl(\prod_{k=j+1}^{t}\alpha_k\Bigr)}_{=:\,a_{t\leftarrow j}}
     \gamma_j\,\bm{B}_j\bm{x}_j,
\end{equation}
where by convention an empty product equals $1$ (so
$a_{t\leftarrow t} = 1$).

\emph{Proof of the induction step.} Assume \eqref{eq:un-claim} at
time $t-1$. Substitute into \eqref{eq:rec2}:
\begin{align}
  \bm{h}_t
   &= \alpha_t\,\bm{h}_{t-1} + \gamma_t\,\bm{B}_t\bm{x}_t \\
   &= \alpha_t \sum_{j=1}^{t-1}
        \Bigl(\prod_{k=j+1}^{t-1}\alpha_k\Bigr)\gamma_j\bm{B}_j\bm{x}_j
        + \gamma_t\,\bm{B}_t\bm{x}_t \\
   &= \sum_{j=1}^{t-1}
        \Bigl(\prod_{k=j+1}^{t}\alpha_k\Bigr)\gamma_j\bm{B}_j\bm{x}_j
        + \gamma_t\,\bm{B}_t\bm{x}_t
        \qquad\text{(absorb $\alpha_t$ into the product)}\\
   &= \sum_{j=1}^{t}
        \Bigl(\prod_{k=j+1}^{t}\alpha_k\Bigr)\gamma_j\bm{B}_j\bm{x}_j.
\end{align}
The last equality uses the empty-product convention for $j=t$. This
proves \eqref{eq:un-claim}.

\subsection{From state to output}

Plug \eqref{eq:un-claim} into the output equation
$\bm{y}_t = \bm{C}_t^\top\bm{h}_t$:
\begin{align}
  \bm{y}_t
   &= \bm{C}_t^\top \sum_{j=1}^{t}
        a_{t\leftarrow j}\,\gamma_j\,\bm{B}_j\bm{x}_j \\
   &= \sum_{j=1}^{t}
        a_{t\leftarrow j}\,\gamma_j\,
        \underbrace{\bm{C}_t^\top \bm{B}_j}_{\text{scalar}}\;\bm{x}_j
        \qquad\text{(linearity of $\bm{C}_t^\top(\cdot)$)} \label{eq:yout} \\
   &= \sum_{j=1}^{t}
        \underbrace{\bigl(a_{t\leftarrow j}\,\gamma_j\bigr)}_{=:\,L_{tj}}
        \bigl(\bm{C}_t^\top \bm{B}_j\bigr)\,\bm{x}_j. \label{eq:yout-L}
\end{align}
The key observation is that $\bm{x}_j$ is a vector while
$\bm{C}_t^\top\bm{B}_j$ is a scalar; we are free to reorder the
factors.

\subsection{Matrix form: $\bm{Y} = (\bm{L}\odot \bm{C}\bm{B}^\top)\bm{X}$}

Define the lower-triangular matrix
\begin{equation}\label{eq:L2}
  L_{tj} \;=\;
  \begin{cases}
     \gamma_j\displaystyle\prod_{k=j+1}^{t}\alpha_k & j \le t,\\[4pt]
     0 & j > t.
  \end{cases}
\end{equation}
Equivalently,
\begin{equation}\label{eq:Lmat}
  \bm{L} \;=\;
  \begin{bmatrix}
     1    &      &       &      \\
     \alpha_2 & 1 &       &      \\
     \alpha_3\alpha_2 & \alpha_3 & 1 &   \\
     \vdots & & \ddots & \\
     \prod_{k=2}^{T}\alpha_k & \cdots & \alpha_T & 1
  \end{bmatrix}\cdot\mathrm{diag}(\gamma_1,\dots,\gamma_T),
\end{equation}
which is the Mamba-2 mask of \cite{mamba3} eq.~(3). Using
\eqref{eq:yout-L} row by row,
\begin{equation}\label{eq:rowform}
  \bm{y}_t \;=\; \sum_{j=1}^{T} L_{tj}\,\bigl(\bm{C}_t^\top \bm{B}_j\bigr)\,\bm{x}_j
         \;=\; \sum_{j=1}^{T}
                \bigl[\,\bm{L}\odot\bm{C}\bm{B}^\top\,\bigr]_{tj}\,\bm{x}_j.
\end{equation}
(Terms with $j>t$ vanish because $L_{tj}=0$.) Stacking over $t$,
\begin{equation}\label{eq:ssd-derived}
  \boxed{\;\bm{Y} \;=\; \bigl(\bm{L}\odot \bm{C}\bm{B}^\top\bigr)\bm{X}.\;}
\end{equation}
This is exactly \eqref{eq:ssd} (\cite{mamba3} eq.~(2)).

\subsection{Reading the SSD identity as linear attention}

Define the standard attention quantities
\begin{equation}
  \bm{Q}_t := \bm{C}_t,\qquad \bm{K}_j := \bm{B}_j,\qquad \bm{V}_j := \bm{x}_j.
\end{equation}
Then \eqref{eq:ssd-derived} is
\begin{equation}\label{eq:attn}
  \bm{Y} \;=\; \bigl(\bm{L}\odot \bm{Q}\bm{K}^\top\bigr)\bm{V},
\end{equation}
which is \emph{linear attention} (no softmax, no row normalization),
with a data-dependent, structured mask $\bm{L}$. Compare to standard
softmax self-attention
$\bm{Y} = \mathrm{softmax}(\bm{Q}\bm{K}^\top / \sqrt{d})\,\bm{V}$:
\begin{itemize}
  \item Standard attention computes a dense $T\times T$ similarity
        matrix $\bm{Q}\bm{K}^\top$ and row-normalizes it by softmax.
  \item Mamba-3 SSD computes the same $\bm{Q}\bm{K}^\top$, then
        multiplies it elementwise by a learned mask $\bm{L}$ whose
        off-diagonal entries decay geometrically in $|t-j|$.
  \item The decay factor $\prod_{k=j+1}^{t}\alpha_k$ plays the role of
        a relative-position bias: tokens separated by large
        $\sum_{k}\Delta_k |A_k|$ contribute exponentially less.
\end{itemize}

\subsection{Trapezoidal three-term form}\label{sec:trap}

The same derivation goes through for \eqref{eq:rec3}. We show only
the new piece: the mask acquires an extra two-band contribution.
Substitute \eqref{eq:rec3} once,
\begin{align}
  \bm{h}_t
   &= \alpha_t\,\bm{h}_{t-1}
     + \beta_t\,\bm{B}_{t-1}\bm{x}_{t-1}
     + \gamma_t\,\bm{B}_t\bm{x}_t.
\end{align}
Iterating (compare \eqref{eq:un-claim}) gives
\begin{equation}\label{eq:un3}
  \bm{h}_t = \sum_{j=1}^{t}
     \Bigl(\prod_{k=j+1}^{t}\alpha_k\Bigr)\!\bigl(\gamma_j\,\bm{B}_j\bm{x}_j
                                               + \beta_j\,\bm{B}_{j-1}\bm{x}_{j-1}\bigr),
\end{equation}
where the $\beta_j$ term contributes to the state via
$\bm{B}_{j-1}\bm{x}_{j-1}$. Plugging into
$\bm{y}_t = \bm{C}_t^\top\bm{h}_t$ and collecting coefficients of
$\bm{B}_j\bm{x}_j$ (note the index shift for the $\beta$ term), one
obtains
\begin{equation}\label{eq:L3}
  L_{tj} \;=\; \Bigl(\prod_{k=j+1}^{t}\alpha_k\Bigr)\gamma_j
            \;+\; \Bigl(\prod_{k=j+2}^{t}\alpha_k\Bigr)\beta_{j+1}
            \quad \text{for } j\le t,
\end{equation}
with the usual empty-product convention. In matrix form this is
\cite{mamba3} eq.~(7): $\bm{L}$ is the product of a
1-semiseparable matrix (whose $(t,j)$ entry is
$\prod_{k=j+1}^{t}\alpha_k$) and a two-band matrix
$\mathrm{diag}(\gamma_0,\dots,\gamma_T) + \mathrm{subdiag}(\beta_1,\dots,\beta_T)$.
Equations \eqref{eq:rowform}–\eqref{eq:ssd-derived} are unchanged.

\paragraph{Boundary conditions of the $\beta$ band.} The second term in
\eqref{eq:L3} references $\beta_{j+1}$. For the left endpoint, the
recurrence \eqref{eq:rec3} at $t=1$ uses $\beta_1 \bm{B}_0\bm{x}_0$
with $\bm{B}_0, \bm{x}_0$ out of range; we therefore adopt the
convention $\beta_1 := 0$, or equivalently drop this term when
$t=1$. Consistently, $\beta_1$ never appears in \eqref{eq:L3}: the
second term at $j=0$ would need $\beta_1$, but $j\ge 1$ by indexing.
For the right endpoint, at $j=t$ the expression $\beta_{j+1}=\beta_{t+1}$
is undefined; we set $\beta_{T+1} := 0$ so the second term vanishes
at $j=t$, leaving $L_{tt} = \gamma_t$ as expected. Any implementation
that pads $\beta$ with zeros on both ends and builds the second band
only for the strictly lower triangle ($j<t$) honours both conventions
automatically. An off-by-one in this shift is silently masked by the
test $\lambda_t\equiv 1$ (which zeroes $\beta$ identically) and must
therefore be caught by an independent test with $\lambda_t \in (0,1)$.

%% ================================================================
\section{Self-Attention via Mamba-3}
\label{sec:self}
%% ================================================================

We now instantiate the duality \eqref{eq:ssd-derived} as a concrete
self-attention block for a single sequence
$\bm{X}\in\mathbb{R}^{T\times D}$.

\paragraph{Why this is essentially a relabelling exercise.}
The SSD identity \eqref{eq:ssd-derived} already has the shape of
attention. Place the two formulas side by side:
\[
  \underbrace{\bm{Y} = \bigl(\bm{L}\odot\bm{C}\bm{B}^\top\bigr)\bm{V}}_{\text{Mamba SSD (2 or 3)}}
  \qquad\Longleftrightarrow\qquad
  \underbrace{\bm{Y} = \mathrm{softmax}\!\bigl(\bm{Q}\bm{K}^\top/\sqrt{d}\bigr)\bm{V}}_{\text{softmax attention}}.
\]
Reading off the substitution that matches the shapes,
\[
  \bm{Q} \;\leftarrow\; \bm{C},\qquad
  \bm{K} \;\leftarrow\; \bm{B},\qquad
  \bm{V} \;\leftarrow\; \mathrm{SiLU}(\bm{W}_V\bm{x}),
\]
i.e.\ the \emph{query} projection becomes $\bm{C}$ and the \emph{key}
projection becomes $\bm{B}$. The value projection $\bm{V} =
\mathrm{SiLU}(\bm{W}_V\bm{x})$ is the same learned $\bm{W}_V$ that
softmax attention uses; it is part of the SSD-as-attention reading
generally (the $\bm{X}$ in the raw identity \eqref{eq:ssd-derived} is
the input, and any attention block needs a learned value), and so is
not specific to Mamba-3 — Mamba-2 SSD-as-attention takes the same
shape \eqref{eq:rec2}–\eqref{eq:ssd-derived} with $\bm{V}$ in place of
$\bm{X}$. The only structural difference from softmax attention is
that row normalisation $\mathrm{softmax}(\cdot/\sqrt{d})$ has been
replaced by an element-wise multiplication by the learned, structured
mask $\bm{L}$. That is the whole operator: pick the right
$\bm{Q}/\bm{K}/\bm{V}$, build $\bm{L}$ from the per-token scalars,
multiply.

\paragraph{The trapezoidal three-term form does not change this shape.}
A reasonable worry: Mamba-3's three-term recurrence \eqref{eq:rec3}
adds a $\beta_t$ band that mixes adjacent tokens, so does the output
formula still factor as $\bm{Y} = (\bm{L}\odot\bm{C}\bm{B}^\top)\bm{V}$?
It does — only the \emph{contents} of $\bm{L}$ change.
\S\ref{sec:trap} derived this already: collecting coefficients of
$\bm{B}_j\bm{x}_j$ in the unrolled three-term state gives
\eqref{eq:L3},
\[
  L_{tj} = \Bigl(\prod_{k=j+1}^{t}\alpha_k\Bigr)\gamma_j
  \;+\; \Bigl(\prod_{k=j+2}^{t}\alpha_k\Bigr)\beta_{j+1},
\]
i.e.\ \emph{two} contributions per $(t,j)$ cell instead of one. The
Mamba-2 two-term mask \eqref{eq:L2} is exactly the special case
$\beta\equiv 0$ (equivalently $\lambda_t\equiv 1$). Plugging this
richer $L_{tj}$ into the same row formula \eqref{eq:rowform}
reproduces $\bm{Y} = (\bm{L}\odot\bm{C}\bm{B}^\top)\bm{V}$ unchanged;
only the mask gets an extra band. Consequence for an implementer:
Mamba-3 self-attention reuses every line of code that handles
Mamba-2 SSD, with the single change of swapping the mask builder
from a two-term to a three-term constructor.

\subsection{Per-token projections}

For each $t\in\{1,\dots,T\}$ compute
\begin{align}
  \bm{B}_t &= \bm{W}_B\,\bm{x}_t        & &\text{(key)} \\
  \bm{C}_t &= \bm{W}_C\,\bm{x}_t        & &\text{(query)} \\
  \Delta_t &= \mathrm{Softplus}(\bm{w}_\Delta^\top\bm{x}_t) & &\text{(step size)}\\
  A_t &= -\mathrm{Softplus}(\bm{w}_A^\top\bm{x}_t) \text{ or const.} & &\text{(decay)} \\
  \lambda_t &= \sigma(\bm{w}_\lambda^\top\bm{x}_t)\in[0,1] & &\text{(trapezoidal weight)} \\
  \bm{v}_t &= \mathrm{SiLU}(\bm{W}_V\,\bm{x}_t)         & &\text{(value)}
\end{align}
Form $\alpha_t,\beta_t,\gamma_t$ from \eqref{eq:abc}. Optionally apply
\textbf{BCNorm} (RMSNorm on $\bm{B}_t,\bm{C}_t$; \cite{mamba3}
\S 3.4) and learnable channel-wise biases $\bm{b}_B, \bm{b}_C$ before
using them.

\subsection{Self-attention output}

Assemble $\bm{B},\bm{C}\in\mathbb{R}^{T\times N}$ and
$\bm{V}\in\mathbb{R}^{T\times D}$ (rows $\bm{v}_t$) and the mask
$\bm{L}$ from \eqref{eq:L2} or \eqref{eq:L3}. The self-attention
output is
\begin{equation}\label{eq:self-out}
  \bm{Y}_{\text{self}} \;=\; \bigl(\bm{L}\odot\bm{C}\bm{B}^\top\bigr)\bm{V},
\end{equation}
or, entrywise,
\begin{equation}\label{eq:self-entry}
  \bm{y}_i^{\text{self}} \;=\; \sum_{j=1}^{i} L_{ij}\,(\bm{C}_i^\top \bm{B}_j)\,\bm{v}_j.
\end{equation}
Because $L_{ij}=0$ for $j>i$, the attention is \emph{causal}: token
$i$ sees tokens $1,\dots,i$. We remove this causality in
Section~\ref{sec:images} by running a second pass in the reverse
direction.

\subsection{Complex-valued SSM / RoPE trick: attention with learned RoPE}
\label{sec:rope}

Mamba-3 promotes $\bm{A}$ from a real scalar to a block-diagonal
matrix of $2{\times}2$ rotations
\begin{equation}\label{eq:Rdef}
  \bm{R}_t = \mathrm{Block}\bigl(\{R(\Delta_t\bm{\theta}_t[i])\}_{i=1}^{N/2}\bigr),
  \qquad
  R(\phi) = \begin{bmatrix}\cos\phi & -\sin\phi\\ \sin\phi & \cos\phi\end{bmatrix},
\end{equation}
and Proposition~4 of \cite{mamba3} shows that this is equivalent
to applying data-dependent rotary embeddings to $\bm{B}$ and $\bm{C}$.
Writing $\bm{R}_{\le t} := \prod_{i=0}^{t}\bm{R}_i^\top$, the complex
recurrence is
\begin{equation}\label{eq:complex-rec}
  \bm{h}_t = \alpha_t\,\bm{h}_{t-1}
            + \beta_t\,\bm{R}_{\le t-1}\,\bm{B}_{t-1}\,\bm{x}_{t-1}
            + \gamma_t\,\bm{R}_{\le t}\,\bm{B}_{t}\,\bm{x}_{t},
\qquad
  \bm{y}_t = (\bm{R}_{\le t}\,\bm{C}_t)^\top\,\bm{h}_t.
\end{equation}
Following the same unrolling as in Section~\ref{sec:ssd}, the
analogue of \eqref{eq:self-entry} becomes
\begin{equation}\label{eq:self-rope}
  \bm{y}_i^{\text{self}}
     \;=\; \sum_{j=1}^{i} L_{ij}\,
            \bigl(\widetilde{\bm{C}}_i^{\,\top}\widetilde{\bm{B}}_j\bigr)\,\bm{v}_j,
  \qquad
  \widetilde{\bm{C}}_i := \bm{R}_{\le i}\,\bm{C}_i,\quad
  \widetilde{\bm{B}}_j := \bm{R}_{\le j}\,\bm{B}_j.
\end{equation}
That is, \textbf{Mamba-3 self-attention is linear attention with a
structured mask $\bm{L}$ and a \emph{data-dependent} rotary embedding}
applied to the $\bm{B}$ (key) and $\bm{C}$ (query) projections,
exactly the Transformer's RoPE but with angles learned per token
rather than fixed to $10000^{-2i/N}$.

\subsection{MIMO extension}\label{sec:mimo}

Mamba-3's MIMO variant (\cite{mamba3} \S 3.3, eqs.~(12)–(14))
stacks $R$ parallel SSMs that share the scalar gates $\alpha_t,\Delta_t$
but carry independent $\bm{B}_t^{(j)},\bm{C}_t^{(i)},\bm{x}_t^{(j)}$.
The head-level output is
\begin{equation}
  \bm{y}_t^{(i)} \;=\; \sum_{j=0}^{R-1}\mathrm{SSM}\bigl(\alpha,\Delta,\bm{B}^{(j)},\bm{C}^{(i)},\bm{x}^{(j)}\bigr)_t,
\end{equation}
i.e.\ $R^2$ SISO calls in parallel. Using \eqref{eq:ssd-derived} on
each, the MIMO self-attention output is
\begin{equation}\label{eq:mimo-self}
  \bm{Y}^{(i)} \;=\; \sum_{j=0}^{R-1}
     \bigl(\bm{L}\odot \bm{C}^{(i)}\bm{B}^{(j)\top}\bigr)\bm{X}^{(j)},
  \quad i=0,\dots,R-1.
\end{equation}
This is exactly multi-head linear attention with shared mask.

%% ================================================================
\section{Cross-Attention via Mamba-3}
\label{sec:cross}
%% ================================================================

For cross-attention we have two token sequences: a \emph{query}
sequence $\bm{X}^{q}\in\mathbb{R}^{T_q\times D}$ and a \emph{key/value}
(context) sequence $\bm{X}^{kv}\in\mathbb{R}^{T_{kv}\times D}$.

\paragraph{Why cross-attention needs no new derivation.}
Cross-attention reuses the same identity, just with two different
sequences feeding the two sides instead of one. Place the two
formulas side by side:
\[
  \underbrace{\bm{Y} = \bigl(\bm{L}^{\text{cross}}\odot\bm{C}^q\bm{B}^{kv\top}\bigr)\bm{V}^{kv}}_{\text{Mamba-3 cross-SSD}}
  \qquad\Longleftrightarrow\qquad
  \underbrace{\bm{Y} = \mathrm{softmax}\!\bigl(\bm{Q}\bm{K}^\top/\sqrt{d}\bigr)\bm{V}}_{\text{softmax cross-attention}},
\]
with the substitution
\[
  \bm{Q} \;\leftarrow\; \bm{C}^q, \quad
  \bm{K} \;\leftarrow\; \bm{B}^{kv}, \quad
  \bm{V} \;\leftarrow\; \mathrm{SiLU}(\bm{W}_V \bm{x}^{kv}).
\]
The reading rule: the \emph{query} projection comes from the
query-side tokens $\bm{x}^{q}$, while \emph{both} the key and the
value projections come from the kv-side tokens $\bm{x}^{kv}$.
Nothing about the operator changes — only which sequence each
projection reads from.

\paragraph{Two things that do change relative to self-attention.}
First, the mask $\bm{L}^{\text{cross}}$ \eqref{eq:Lcross} is
rectangular ($T_q \times T_{kv}$) rather than triangular, because
queries and keys live on different time axes and there is no notion
of one token coming "before" another across the two streams.
Second, the trapezoidal structure of \eqref{eq:L3} carries over to
the kv axis when the kv side is itself produced by a Mamba-3
recurrence (the typical case in Variant~B, \S\ref{sec:crossB}): each
column of $\bm{L}^{\text{cross}}$ then has the same two-band decay
profile as a self-SSD row, with the recency bias discussed in
\S\ref{sec:crossB}. The output formula $\bm{Y} =
(\bm{L}^{\text{cross}}\odot\bm{C}^q\bm{B}^{kv\top})\bm{V}^{kv}$ is
\emph{unchanged}; only the geometry of $\bm{L}$ and which tokens
supply $\bm{C}, \bm{B}, \bm{V}$ are different. Consequence for an
implementer: cross-attention reuses every line of code that handles
self-SSD, with two changes — feed $\bm{C}$ from the query sequence
and $\bm{B}, \bm{V}$ from the kv sequence, and swap the
triangular mask builder for the rectangular \texttt{cross} mask.

\subsection{Projections for the two sides}

Query side:
\begin{equation}\label{eq:cross-proj-q}
  \bm{C}_i^{q} = \bm{W}_C\,\bm{x}_i^{q},
  \quad i=1,\dots,T_q,
\end{equation}
that is, only the \emph{query} projection is taken from the query
side. Key/value side:
\begin{align}
  \bm{B}_j^{kv} &= \bm{W}_B\,\bm{x}_j^{kv}, \\
  \Delta_j^{kv} &= \mathrm{Softplus}(\bm{w}_\Delta^\top \bm{x}_j^{kv}), \\
  A_j^{kv}      &= -\mathrm{Softplus}(\bm{w}_A^\top \bm{x}_j^{kv})\;\text{or const.}, \\
  \bm{v}_j^{kv} &= \mathrm{SiLU}(\bm{W}_V\,\bm{x}_j^{kv}),
  \quad j=1,\dots,T_{kv}.
\end{align}
Decay scalars $\alpha_j^{kv}$ etc.\ are computed as in \eqref{eq:abc}
but using the $kv$-side $\Delta_j^{kv}, A_j^{kv}$.

\subsection{Variant A: state-compressed cross-attention (constant memory)}
\label{sec:crossA}

Run the Mamba-3 recurrence \eqref{eq:rec2} on the $kv$-side with
\emph{outer-product} accumulation, treating $\bm{v}_j^{kv}$ as the
input instead of $\bm{x}_j^{kv}$. Define a matrix-valued hidden state
$\bm{H}_j \in \mathbb{R}^{N\times D}$ by
\begin{equation}\label{eq:Hrec}
  \bm{H}_j \;=\; \alpha_j^{kv}\,\bm{H}_{j-1}
                + \gamma_j^{kv}\,\bm{B}_j^{kv}\,\bm{v}_j^{kv\,\top},
  \qquad \bm{H}_0 = \bm{0}.
\end{equation}
By the exact same unrolling as in \eqref{eq:un-claim} we obtain
\begin{equation}\label{eq:Hclose}
  \bm{H}_{T_{kv}} \;=\; \sum_{j=1}^{T_{kv}}
      \Bigl(\prod_{k=j+1}^{T_{kv}}\alpha_k^{kv}\Bigr)
      \gamma_j^{kv}\,\bm{B}_j^{kv}\,\bm{v}_j^{kv\,\top}
  \;=:\; \bm{h}_{\text{ref}} \in \mathbb{R}^{N\times D}.
\end{equation}
This $\bm{h}_{\text{ref}}$ is the \emph{compressed context}: it
condenses the entire $kv$-sequence into a single $N\times D$ matrix.
For each query $\bm{C}_i^q$, the output is obtained by a single
left-multiplication:
\begin{equation}\label{eq:crossA-out}
  \boxed{\;\bm{y}_i^{\text{cross-A}} \;=\; \bm{C}_i^{q\,\top}\,\bm{h}_{\text{ref}} \;\in \mathbb{R}^{D}.\;}
\end{equation}
Shape check: $\bm{C}_i^{q\,\top}\in\mathbb{R}^{1\times N}$ times
$\bm{h}_{\text{ref}}\in\mathbb{R}^{N\times D}$ gives
$\mathbb{R}^{1\times D}$, i.e.\ a $D$-vector.

\textbf{Cost.} Once $\bm{h}_{\text{ref}}$ is built in $O(T_{kv}ND)$
time, answering all $T_q$ queries costs $O(T_q ND)$. Memory for the
``KV cache'' is constant in $T_{kv}$ (it is exactly
$\bm{h}_{\text{ref}}\in\mathbb{R}^{N\times D}$). This is the
analogue of the recurrent KV cache.

\subsection{Variant B: token-level cross-attention (exact SSD)}
\label{sec:crossB}

To allow every query token to see every $kv$ token independently,
keep the full similarity matrix. Build an $T_q\times T_{kv}$ mask
\begin{equation}\label{eq:Lcross}
  L^{\text{cross}}_{ij} \;=\; \gamma_j^{kv}\,
        \Bigl(\prod_{k=j+1}^{T_{kv}}\alpha_k^{kv}\Bigr),
\end{equation}
whose $j$-th column is the same decay profile as in \eqref{eq:L2}, but
now the matrix is \emph{not triangular} because queries do not live on
the same time axis as keys. The cross-attention output is
\begin{equation}\label{eq:crossB-out}
  \boxed{\;\bm{Y}^{\text{cross-B}} \;=\; \bigl(\bm{L}^{\text{cross}}\odot\bm{C}^{q}\bm{B}^{kv\,\top}\bigr)\bm{V}^{kv},\;}
\end{equation}
or, entrywise,
\begin{equation}\label{eq:crossB-entry}
  \bm{y}_i^{\text{cross-B}} \;=\; \sum_{j=1}^{T_{kv}} L^{\text{cross}}_{ij}\,
           \bigl(\bm{C}_i^{q\,\top}\bm{B}_j^{kv}\bigr)\,\bm{v}_j^{kv}.
\end{equation}
Compare with softmax cross-attention
$\bm{Y} = \mathrm{softmax}(\bm{C}^q\bm{B}^{kv\,\top})\,\bm{V}^{kv}$:
\eqref{eq:crossB-out} replaces softmax row normalization with an
elementwise, structured, learned mask. With the trivial mask
$L^{\text{cross}}_{ij}\equiv 1$ we obtain ordinary linear cross-
attention $\bm{C}^q\bm{B}^{kv\,\top}\bm{V}^{kv}$.

\paragraph{End-of-sequence bias in $\bm{L}^{\text{cross}}$.} Note that
\eqref{eq:Lcross} does not depend on $i$: every query row is the same
column-decay profile. Concretely, $L^{\text{cross}}_{i,T_{kv}}=\gamma_{T_{kv}}$
while $L^{\text{cross}}_{i,1}=\gamma_1\prod_{k=2}^{T_{kv}}\alpha_k$ is
exponentially small once $T_{kv}$ is large. This encodes a strong
\emph{recency bias toward the end of the $kv$ sequence}: every query
already prefers the last $kv$ tokens before any data-dependent
similarity $\bm{C}^q_i{}^\top\bm{B}^{kv}_j$ is applied. This is the
right inductive bias when the $kv$ sequence has a natural causal order
(decoder-style cross-attention), but it is arbitrary when the $kv$
sequence is a permutation-invariant set — most notably, when it comes
from an image patch grid. The bias becomes the order chosen by the
raster scan.

\paragraph{Bidirectional cross-attention for permutation-invariant $kv$.}
When the $kv$ stream has no preferred order, replace
$\bm{L}^{\text{cross}}$ by the symmetrised mask
\begin{equation}\label{eq:Lcross-bi}
  L^{\text{cross},\leftrightarrow}_{ij}
     \;=\; \gamma_j\Bigl[\prod_{k=j+1}^{T_{kv}}\alpha_k
                         \;+\; \prod_{k=1}^{j-1}\alpha_k\Bigr],
\end{equation}
which is obtained by summing \eqref{eq:Lcross} with its reverse-order
counterpart, analogously to \eqref{eq:bidi2}. The two products coincide
at the centre of the $kv$ sequence and decay symmetrically toward both
endpoints, eliminating the raster-order preference while preserving the
linear-complexity structure. For images this is the natural cross-
attention mask to pair with the bidirectional self-attention of
Section~\ref{sec:bidi}.

\subsection{Which variant to use}

\begin{center}
\begin{tabular}{|l|l|l|}
\hline
 & \textbf{Variant A} (state-compressed) & \textbf{Variant B} (token-level)\\ \hline
State size        & $\mathcal{O}(ND)$             & $\mathcal{O}(T_q T_{kv})$ (or streamed) \\ \hline
Compute          & $\mathcal{O}((T_{kv}+T_q)ND)$ & $\mathcal{O}(T_q T_{kv}(N+D))$ \\ \hline
Query adaptivity & Only via $\bm{C}_i^q$         & Via full $\bm{C}_i^{q\top}\bm{B}_j^{kv}$ \\ \hline
Typical use      & Long context / large $T_{kv}$ & Short cross-attention / ablations \\ \hline
\end{tabular}
\end{center}

Both variants can be augmented with the RoPE trick of
\S\ref{sec:rope}: apply the rotation $\bm{R}_{\le i}^{(q)}$ to each
query $\bm{C}_i^q$ and $\bm{R}_{\le j}^{(kv)}$ to each key
$\bm{B}_j^{kv}$ before the inner product. This injects relative
positional information into cross-attention.

%% ================================================================
\section{Applying Mamba-3 to Images}
\label{sec:images}
%% ================================================================

We now assemble the self- and cross-attention blocks into an image
model.

\subsection{Patchification}

Given an image $\bm{I}\in\mathbb{R}^{H\times W\times 3}$, choose a
patch size $P$ and define $L = (H/P)(W/P)$ patches
$\bm{p}_{r,c}\in\mathbb{R}^{3P^2}$ for
$r\in\{1,\dots,H/P\},\,c\in\{1,\dots,W/P\}$. Project with a shared
linear map:
\begin{equation}\label{eq:patch-embed}
  \bm{x}_{r,c} \;=\; \bm{W}_{\text{emb}}\,\bm{p}_{r,c} + \bm{b}_{\text{emb}},
  \qquad \bm{W}_{\text{emb}}\in\mathbb{R}^{D\times 3P^2}.
\end{equation}
Flattening in row-major order gives a token sequence
$\bm{x}_1,\dots,\bm{x}_L$. The same block works for $N_v$ views
stacked as $N_v L$ tokens.

\subsection{Why a plain causal scan fails on images}

Applying \eqref{eq:self-entry} directly to the flattened sequence
gives
\begin{equation}
  \bm{y}_i^{\text{self}} = \sum_{j=1}^{i} L_{ij}\,(\bm{C}_i^\top\bm{B}_j)\,\bm{v}_j,
\end{equation}
i.e.\ each patch only sees patches \emph{earlier in raster order}.
A patch at the top-right thus never reads from the bottom-left — half
the image is missing. An image has no natural causal direction, so
this asymmetry is harmful.

\subsection{Bidirectional and multi-direction scans}
\label{sec:bidi}

Fix a permutation $\pi$ of $\{1,\dots,L\}$. Let
$\mathrm{SSD}_\pi(\bm{X})$ mean: permute rows by $\pi$, apply
\eqref{eq:ssd-derived}, permute back. Because SSD's only source of
order is the mask $\bm{L}$, every permutation gives a different,
still-valid SSM scan.

\paragraph{Two-direction (forward + reverse).} With the identity
permutation $\pi_\to$ and its reverse $\pi_\leftarrow$,
\begin{equation}\label{eq:bidi2}
  \bm{Y}^{\text{bi}}
   \;=\; \mathrm{SSD}_{\pi_\to}(\bm{X}) + \mathrm{SSD}_{\pi_\leftarrow}(\bm{X}).
\end{equation}
Unrolling: for any pair $(i,j)$, either $j\le i$ (covered by the
forward pass) or $j>i$ (covered by the reverse pass), so every pair
of tokens contributes exactly once, with a mask value depending on
$|i-j|$ rather than on the sign of $i-j$. This turns causal SSD into
bidirectional linear attention on the raster order, matching the
bidirectional Mamba variants used in the literature.

\paragraph{Four-direction (raster + reverse + column-major + reverse).}
For 2-D structure, also scan along columns. Let $\pi^\downarrow$ scan
column-by-column top-to-bottom and $\pi^\uparrow$ its reverse. Then
\begin{equation}\label{eq:bidi4}
  \bm{Y}^{\text{4-dir}}
   = \mathrm{SSD}_{\pi_\to}(\bm{X}) + \mathrm{SSD}_{\pi_\leftarrow}(\bm{X})
   + \mathrm{SSD}_{\pi^\downarrow}(\bm{X}) + \mathrm{SSD}_{\pi^\uparrow}(\bm{X}).
\end{equation}
Each summand uses the same learned weights; only the token order
changes, so \eqref{eq:bidi4} adds no parameters. This mirrors the
four scans of Vision-Mamba / V-Mamba realized under the Mamba-3 SSD
kernel.

\paragraph{Equivalence to a symmetric mask.} Writing
\eqref{eq:bidi2} entrywise,
\begin{equation}\label{eq:bidi-sym}
  \bm{y}_i^{\text{bi}} = \sum_{j=1}^{L} \widetilde{L}_{ij}\,
       (\bm{C}_i^\top\bm{B}_j)\,\bm{v}_j,
  \quad
  \widetilde{L}_{ij} = L_{ij}^{\to} + L_{ij}^{\leftarrow},
\end{equation}
where $L^\to$ is \eqref{eq:L2} and $L^\leftarrow$ is the same matrix
built on the reversed sequence. $\widetilde{L}$ is \emph{symmetric} in
$i,j$: the forward pass covers $j\le i$, the reverse pass covers
$j\ge i$, and with the diagonal counted twice the bidirectional
self-attention has no preferred direction.

\subsection{2-D positional encoding via the complex SSM}
\label{sec:rope2d}

To encode row and column positions, use the RoPE trick of
\S\ref{sec:rope} with a \emph{factored} angle schedule. For the
$n$-th channel of the state, split into independent row and column
angles
\begin{equation}\label{eq:rope2d}
  \theta_{r,c}[n] = \theta^{\text{row}}[n]\cdot r + \theta^{\text{col}}[n]\cdot c,
  \qquad n = 1,\dots,N/2,
\end{equation}
where $\theta^{\text{row}},\theta^{\text{col}}$ are learned or fixed
frequency vectors. Replace $\bm{\theta}_t$ in \eqref{eq:Rdef} by
$\theta_{r,c}$ for the patch at $(r,c)$. The resulting
$\widetilde{\bm{C}}_i = \bm{R}_{\le i}\bm{C}_i$ and
$\widetilde{\bm{B}}_j$ now carry a 2-D rotary position encoding
inside the SSD inner product \eqref{eq:self-rope}.

\subsection{Image self-attention layer}

Putting the pieces together, one Mamba-3 self-attention layer for
images is:

\medskip
\noindent\textbf{Algorithm (Mamba-3 image self-attention layer).}
\begin{enumerate}
  \item Input tokens $\bm{X}\in\mathbb{R}^{L\times D}$ for one image
        (already patch-embedded).
  \item Compute per-token $\bm{B}_t,\bm{C}_t,\Delta_t,A_t,\lambda_t,\bm{v}_t$
        via \eqref{eq:projB}–\eqref{eq:projV}.
  \item Apply BCNorm and add learnable biases $\bm{b}_B,\bm{b}_C$
        (\cite{mamba3}~\S 3.4).
  \item Apply 2-D RoPE: $\bm{B}_t\leftarrow \bm{R}_{\le t}\bm{B}_t$,
        $\bm{C}_t\leftarrow \bm{R}_{\le t}\bm{C}_t$
        with $\bm{R}$ from \eqref{eq:rope2d}.
  \item Build masks $\bm{L}^{\to},\bm{L}^{\leftarrow}$ from
        \eqref{eq:L2} (or \eqref{eq:L3} for the trapezoidal
        variant) using $\alpha_t,\beta_t,\gamma_t$.
  \item Compute $\bm{Y}^\to = (\bm{L}^\to\odot\bm{C}\bm{B}^\top)\bm{V}$
        and $\bm{Y}^\leftarrow$ on the reversed order; optionally
        also $\bm{Y}^\downarrow,\bm{Y}^\uparrow$.
  \item Sum across scans to get
        $\bm{Y}^{\text{self}}\in\mathbb{R}^{L\times D}$.
  \item Apply the usual post-processing (output projection +
        SwiGLU MLP block + residual), exactly as in a Mamba-3 /
        Llama block (\cite{mamba3}~Fig.~2).
\end{enumerate}

\subsection{Multi-view cross-attention for 3-D applications}
\label{sec:multiview}

For multi-view inputs (e.g.\ \cite{depthanything3}, Fig.~2),
stack the tokens of $N_v$ views into
$\bm{X}\in\mathbb{R}^{N_v L \times D}$. Let $V(i)\in\{1,\dots,N_v\}$
denote the view index of token $i$.

\paragraph{Within-view layer.} Apply the self-attention layer of the
previous subsection \emph{per view}: i.e.\ the mask is
$L_{ij}^{\text{in}} = L_{ij}\cdot\mathbb{I}[V(i)=V(j)]$, so tokens
only interact inside their own view. This is block-diagonal
bidirectional SSD.

\paragraph{Cross-view layer.} Treat one view at a time as the
\emph{query} side and all other views as the \emph{key/value} side.
Build
\begin{equation}\label{eq:Lxview}
  L_{ij}^{\text{cv}}
    \;=\; L^{\text{cross}}_{ij}\cdot \mathbb{I}[V(i)\neq V(j)],
\end{equation}
with $L^{\text{cross}}$ from \eqref{eq:Lcross} (where the ``time'' axis
is a learned ordering of views, e.g.\ a canonical ``reference view
first'' order). The output is
\begin{equation}\label{eq:cvout}
  \bm{Y}^{\text{cv}} \;=\; \bigl(\bm{L}^{\text{cv}}\odot\bm{C}\bm{B}^\top\bigr)\bm{V}.
\end{equation}
Alternating within-view and cross-view layers realizes the same
design as \cite{depthanything3} Fig.~2 (``within-view
self-attention'' and ``cross-view self-attention''), but with softmax
attention replaced by the Mamba-3 SSD kernel, inheriting its linear
complexity in the token count.

\paragraph{Two-view cross-attention (asymmetric).} For the special
case of a query image and a reference image, use Variant~A of
Section~\ref{sec:crossA}: build
$\bm{h}_{\text{ref}}\in\mathbb{R}^{N\times D}$ from the reference
tokens once, and answer each query token in $O(ND)$. This is the
natural ``Mamba-3 cross-attention'' block when the reference view is
much larger than the query view, or is reused across many queries
(e.g.\ a novel-view-synthesis decoder reading from a shared scene
encoding).

%% ================================================================
\section{Non-Causal SSD and the VSSD Image Model}
\label{sec:ncssd}
%% ================================================================

The bidirectional and multi-direction scans of \S\ref{sec:bidi}
remove causality by \emph{repeating} the SSD kernel along several
permutations of the token order. \cite{vssd} (VSSD, ICCV~2025)
proposes an alternative that removes causality \emph{at the algebraic
level}, so a single forward pass already yields a non-causal,
position-independent operator. We derive their construction from the
two-term recurrence \eqref{eq:rec2} using only the algebra of
Section~\ref{sec:ssd}, then describe the resulting image block.

\subsection{The reinterpretation of $A_t$}

Recall \eqref{eq:rec2}, $\bm{h}_t = \alpha_t\,\bm{h}_{t-1}
+ \gamma_t\,\bm{B}_t\bm{x}_t$ with $\alpha_t = e^{\Delta_t A_t}\in
(0,1)$ controlling \emph{how much of the previous state is retained}.
\cite{vssd} observe that in the SSD form, $A_t$ is a learned,
data-dependent scalar with no a priori meaning: one is equally free to
make it modulate the \emph{current} contribution instead. Concretely,
fix the previous-state multiplier to $1$ and divide the input weight
by $A_t$:
\begin{equation}\label{eq:nc-rec}
  \bm{h}_t \;=\; \bm{h}_{t-1} \;+\; \tfrac{1}{A_t}\,\bm{B}_t\bm{x}_t,
  \qquad \bm{h}_0 = \bm{0},
\end{equation}
which is \cite{vssd}~eq.~(7). Unrolling \eqref{eq:nc-rec} with
$\bm{h}_0=\bm{0}$,
\begin{equation}\label{eq:nc-unroll}
  \bm{h}_t \;=\; \sum_{i=1}^{t} \tfrac{1}{A_i}\,\bm{B}_i\bm{x}_i.
\end{equation}
Compared with \eqref{eq:un-claim}, the multiplicative decay
$\prod_{k=j+1}^{t}\alpha_k$ has been replaced by a flat $1$: a token's
contribution to the state no longer depends on \emph{when} it
arrives, only on its own coefficient $1/A_i$. The mechanism is now
\emph{self-referential}: each token decides its own weight.

\subsection{Bidirectional NC-SSD collapses to a global hidden state}

The recurrence \eqref{eq:nc-unroll} is still causal in the sense that
$\bm{h}_t$ sums only over $i\le t$. Run it once forward and once
backward and sum, exactly as in \eqref{eq:bidi2}, defining
$\bm{Z}_j := \bm{B}_j\bm{x}_j^\top \in \mathbb{R}^{N\times D}$ (the
outer product appearing in \eqref{eq:Hrec}):
\begin{align}
  \bm{H}_i^{\,\to}        &= \sum_{j=1}^{i}\tfrac{1}{A_j}\bm{Z}_j, \\
  \bm{H}_i^{\,\leftarrow} &= \sum_{j=i}^{L}\tfrac{1}{A_j}\bm{Z}_j, \\
  \bm{H}_i^{\,\text{bi}}  &= \bm{H}_i^{\,\to} + \bm{H}_i^{\,\leftarrow} - \tfrac{1}{A_i}\bm{Z}_i
                          \;=\; \sum_{j=1}^{L}\tfrac{1}{A_j}\bm{Z}_j. \label{eq:nc-bi}
\end{align}
The $-\tfrac{1}{A_i}\bm{Z}_i$ correction removes the double counting
of $j=i$. The last equality is the key observation of
\cite{vssd}: the bidirectional NC-SSD hidden state is
\emph{independent of $i$}. Writing
\begin{equation}\label{eq:nc-H}
  \bm{H} \;:=\; \sum_{j=1}^{L}\tfrac{1}{A_j}\bm{Z}_j
       \;=\; \sum_{j=1}^{L}\tfrac{1}{A_j}\,\bm{B}_j\,\bm{x}_j^\top
       \;\in\;\mathbb{R}^{N\times D},
\end{equation}
every output token reads from the \emph{same} compressed context
$\bm{H}$. \cite{vssd}~eq.~(8) further notes that the self-term
$\tfrac{1}{A_i}\bm{Z}_i$ is redundant with the block's skip
connection, so in practice one keeps the simpler global form
\eqref{eq:nc-H} and skips the correction.

\subsection{Closed form: $\bm{Y} = \bm{C}\bigl((\bm{B}\odot\bm{m})^\top\bm{X}\bigr)$}

Let $\bm{m}\in\mathbb{R}^{L}$ with $m_j := 1/A_j$, and write $\bm{B}\odot\bm{m}$
for the matrix whose $j$-th row is $m_j\,\bm{B}_j$ (i.e.\ $\bm{m}$
broadcast along the $N$-axis). The output equation
$\bm{y}_t = \bm{C}_t^\top\bm{H}$ stacks to
\begin{equation}\label{eq:nc-Y}
  \boxed{\;\bm{Y} \;=\; \bm{C}\,\bm{H} \;=\; \bm{C}\bigl((\bm{B}\odot\bm{m})^\top\bm{X}\bigr),\;}
\end{equation}
which is \cite{vssd}~eq.~(10). Shape check:
$(\bm{B}\odot\bm{m})^\top\in\mathbb{R}^{N\times L}$, $\bm{X}\in
\mathbb{R}^{L\times D}$, product is $\mathbb{R}^{N\times D}$, and
$\bm{C}\in\mathbb{R}^{L\times N}$ times that is
$\mathbb{R}^{L\times D}$, as required.

Comparing \eqref{eq:nc-Y} with the SSD identity
\eqref{eq:ssd-derived},
\begin{equation}
  \bm{Y}_{\text{SSD}} \;=\; \bigl(\bm{L}\odot\bm{C}\bm{B}^\top\bigr)\bm{X},
  \qquad
  \bm{Y}_{\text{NC-SSD}} \;=\; \bm{C}\bigl((\bm{B}\odot\bm{m})^\top\bm{X}\bigr),
\end{equation}
the $T\times T$ structured mask $\bm{L}$ has been replaced by a
$T$-vector $\bm{m}$ that column-scales $\bm{B}$. The two operations
are no longer of the same shape:
\begin{itemize}
  \item SSD pays $O(T^2 N)$ to form $\bm{L}\odot\bm{C}\bm{B}^\top$,
        then $O(T^2 D)$ to multiply by $\bm{X}$, with state of size
        $\bm{L}$ ($\Theta(T^2)$ if materialized).
  \item NC-SSD pays $O(TND)$ to form $\bm{H}$, then $O(TND)$ for
        $\bm{C}\bm{H}$, with state of size $\bm{H}\in\mathbb{R}^{N\times D}$
        ($O(ND)$, \emph{independent of $T$}).
\end{itemize}
The price is uniformity: every query reads the same global summary
$\bm{H}$, so the per-pair relative weight that $\bm{L}_{ij}$ encodes
in SSD is gone. \cite{vssd} argues this is a correct trade for
images, where no $1$-D distance is meaningful anyway.

\paragraph{Practical note on parameterization.} \cite{vssd}~\S 3.2
remarks that since $A_t$ and $1/A_t$ span the same range of values up
to inversion, it is numerically simpler to learn $\bm{m}$ directly
(equivalently, to learn $A$ via the existing projection
\eqref{eq:projA} and then treat its output as $\bm{m}$). The two
parameterizations differ only in the optimization geometry.

\subsection{Relation to Mamba-3}\label{sec:ncssd-mamba3}

VSSD is built on Mamba-2 (scalar $A$, two-term recurrence). The
non-causality argument carries over to Mamba-3 unchanged: at the SSD
level, both Mamba-2's \eqref{eq:L2} and Mamba-3's \eqref{eq:L3} have
the same shape $\bm{Y} = (\bm{L}\odot\bm{C}\bm{B}^\top)\bm{X}$, and
setting $\alpha_t \equiv 1$, $\gamma_t \mapsto 1/A_t$, $\beta_t \equiv 0$
turns \eqref{eq:L3} into the rank-one mask $\bm{L}_{ij} = m_j$, which
factors out of the Hadamard product:
\begin{equation}
  \bm{L}_{ij} = m_j
  \;\Longrightarrow\;
  \bm{L}\odot\bm{C}\bm{B}^\top
     = \bm{C}\,(\bm{B}\odot\bm{m})^\top,
\end{equation}
recovering \eqref{eq:nc-Y}. The Mamba-3 extras compose cleanly on top
of NC-SSD:
\begin{itemize}
  \item The RoPE trick of \S\ref{sec:rope} still applies: rotate
        $\bm{B}_j, \bm{C}_i$ by $\bm{R}_{\le j}, \bm{R}_{\le i}$ before
        \eqref{eq:nc-Y}, and the inner product
        $\widetilde{\bm{C}}_i^\top \widetilde{\bm{B}}_j$ carries relative
        positional information. With the 2-D factorization
        \eqref{eq:rope2d}, this becomes the 2-D RoPE used by VSSD.
  \item The MIMO extension of \S\ref{sec:mimo} carries over by
        replicating \eqref{eq:nc-Y} per head: $\bm{Y}^{(i)} =
        \bm{C}^{(i)}\bigl((\bm{B}^{(j)}\odot\bm{m})^\top\bm{X}^{(j)}\bigr)$,
        summed over $j$.
\end{itemize}
The trapezoidal three-term form \eqref{eq:rec3} does \emph{not}
survive the NC reduction: $\beta_t$ implements a two-band relative
bias that fundamentally depends on $i,j$ separation, which the
single-vector $\bm{m}$ cannot represent. VSSD-style NC-SSD is thus
strictly a Mamba-2-style reduction of Mamba-3 — useful when the
mask's expressivity is not needed (Section~\ref{sec:ncssd-when}).

\subsection{Tensor-contraction algorithm}\label{sec:ncssd-tensor}

\cite{vssd}~eq.~(9) writes \eqref{eq:nc-Y} as three einsum
contractions, which matches the standard SSD kernel up to one
collapsing step:
\begin{align}
  \bm{Z} &= \mathrm{contract}(LD, LN \to LND)(\bm{X}, \bm{B}), \label{eq:nc-Z}\\
  \bm{H} &= \mathrm{contract}(L, LND \to ND)(\bm{m}, \bm{Z}), \label{eq:nc-H2}\\
  \bm{Y} &= \mathrm{contract}(LN, ND \to LD)(\bm{C}, \bm{H}). \label{eq:nc-Y2}
\end{align}
Step~\eqref{eq:nc-Z} expands each token to its outer product
$\bm{B}_j\bm{x}_j^\top$; step~\eqref{eq:nc-H2} reduces the $L$ axis
into the single global state $\bm{H}$; step~\eqref{eq:nc-Y2} reads
the state out per query. Compared with SSD's
$\mathrm{contract}(LL, LDN \to ND)$ (\cite{vssd}~eq.~(9), top), the
$LL$ mask has degenerated to an $L$-vector and the algorithm has no
explicit $T\times T$ tensor at any point.

\subsection{The VSSD image block}

The NC-SSD operator \eqref{eq:nc-Y} replaces the SSD kernel inside an
otherwise standard vision block. Following \cite{vssd}~\S 3.3:
\begin{enumerate}
  \item \textbf{Local Perception Unit (LPU)} \cite{vssd}: a depthwise
        $3{\times}3$ convolution applied to the patch grid before the
        NC-SSD and again before the FFN, injecting local 2-D
        inductive bias on top of the global, position-independent
        NC-SSD context.
  \item \textbf{Depth-wise conv 1-D $\to$ 2-D.} The causal Conv1D
        used inside Mamba-2 between the input and the SSD kernel is
        replaced by a non-causal depth-wise Conv2D with kernel
        size~$3$, consistent with the model being non-causal.
  \item \textbf{2-D RoPE on $\bm{B},\bm{C}$.} The factored angle
        \eqref{eq:rope2d} is applied to $\bm{C}_i$ and $\bm{B}_j$ before
        the inner product in \eqref{eq:nc-Y}, restoring the spatial
        relationship between patches that is otherwise lost in the
        permutation-invariant global state $\bm{H}$. This is the
        Mamba-3 RoPE trick of \S\ref{sec:rope} applied to the NC-SSD
        kernel.
  \item \textbf{FFN, residual, normalization.} A standard
        Transformer-style FFN follows the NC-SSD; both the NC-SSD and
        the FFN are wrapped in residual connections, matching the
        ViT / Swin block layout rather than Mamba's pre-norm SwiGLU
        layout.
  \item \textbf{Hybrid with MSA in the last stage.} The four-stage
        backbone uses NC-SSD blocks in stages~1--3 and standard
        softmax MSA in the last stage. The rationale is that early
        stages have many tokens, where the $O(T)$ scaling of NC-SSD
        dominates; the last stage has few tokens but high-level
        semantics, where the per-pair adaptivity of softmax MSA pays
        off.
  \item \textbf{Overlapped downsampling} between stages, in the
        style of \cite{vssd}, replaces the non-overlapped patch
        merging of VMamba.
\end{enumerate}
\paragraph{The vector $\bm{m}$ as a saliency map.} Because $\bm{m}\in
\mathbb{R}^{L}$ acts as the only token-dependent weight in
\eqref{eq:nc-Y}, it can be visualized as a 2-D heat map over the patch
grid. Empirically \cite{vssd} report that $\bm{m}$ concentrates on
the foreground object of the image — i.e.\ NC-SSD ends up learning a
saliency-like weighting that the SSD mask $\bm{L}$ never makes
explicit.

\subsection{When to use NC-SSD versus bidirectional SSD}
\label{sec:ncssd-when}

\begin{center}
\begin{tabular}{|l|l|l|}
\hline
 & \textbf{Bidirectional SSD} (\S\ref{sec:bidi}) & \textbf{NC-SSD} (this section)\\ \hline
Mask state         & $O(T^2)$ (or banded fused kernel) & $O(ND)$ \\ \hline
Compute per layer  & $O(T^2(N+D))\cdot k_{\text{scans}}$ & $O(TND)$ \\ \hline
Per-pair adaptivity & $\bm{L}_{ij}$ and $\bm{C}_i^\top\bm{B}_j$ & $\bm{C}_i^\top\bm{B}_j$ only \\ \hline
Position encoding   & decay in $|i-j|$ + RoPE   & RoPE only (decay collapses) \\ \hline
Typical regime      & moderate $T$, structure-rich tokens & long $T$, vision-style tokens \\ \hline
\end{tabular}
\end{center}
The two are not mutually exclusive: a hybrid stack can use NC-SSD in
the early, high-resolution stages where $T$ dominates the cost, and
bidirectional SSD (or softmax MSA, as in VSSD) in deeper stages where
the mask's per-pair weighting matters. For our multi-view setting
(\S\ref{sec:multiview}), NC-SSD also offers a free state-compressed
cross-view block: each view's NC-SSD hidden state
$\bm{H}^{(v)}\in\mathbb{R}^{N\times D}$ is already in the form that
\eqref{eq:crossA-out} consumes, so cross-view attention reduces to
$\bm{C}^{q\,\top}\sum_v \bm{H}^{(v)}$ with no additional kernel.

%% ================================================================
\section{End-to-End Summary}
%% ================================================================

Using only \eqref{eq:ssd-derived}, we have obtained:

\begin{itemize}
  \item \textbf{Self-attention:} \eqref{eq:self-out} and, with
        learned RoPE, \eqref{eq:self-rope}. Multi-head via
        \eqref{eq:mimo-self}.
  \item \textbf{Cross-attention:} Variant~A \eqref{eq:crossA-out}
        (cheap, compressed KV) and Variant~B \eqref{eq:crossB-out}
        (exact).
  \item \textbf{Image encoder:} patchify + bidirectional /
        multi-direction SSD self-attention
        (\eqref{eq:bidi2},\eqref{eq:bidi4}) with 2-D RoPE
        \eqref{eq:rope2d}.
  \item \textbf{Multi-view encoder:} alternating within-view
        \eqref{eq:cvout}-style blocks and cross-view
        \eqref{eq:Lxview} blocks; or the compressed-state cross
        block \eqref{eq:crossA-out} when one side dominates.
  \item \textbf{Non-causal SSD (VSSD):} reinterpret $A_t$ as a
        current-token weight to obtain \eqref{eq:nc-Y},
        $\bm{Y} = \bm{C}\bigl((\bm{B}\odot\bm{m})^\top\bm{X}\bigr)$,
        which collapses the $T\times T$ mask $\bm{L}$ to a
        $T$-vector $\bm{m}$, removes causality without multi-scan,
        and yields the VSSD image block (Section~\ref{sec:ncssd}).
\end{itemize}

Every output formula above is a direct application of the single
identity
\begin{equation*}
  \bm{Y} = (\bm{L}\odot\bm{C}\bm{B}^\top)\bm{X},
\end{equation*}
which we derived from the Mamba-3 recurrence in
Section~\ref{sec:ssd}. The choices made by a specific architecture
amount to (i) which tokens supply $\bm{B},\bm{C},\bm{X}$ (self vs
cross vs multi-view), and (ii) which structured mask $\bm{L}$ is used
(causal vs bidirectional vs block-diagonal vs view-masked). Mamba-3's
exponential-trapezoidal rule \eqref{eq:L3} and RoPE trick
\eqref{eq:complex-rec} then give the mask and the query/key its final,
expressive form.

\bibliographystyle{plain}
\bibliography{references}

\end{document}
